\section{The Motivating Attack: Minting for a Guess}
\label{sec:guess-to-mint}

The protocol of Sections~\ref{sec:protocol}--\ref{sec:verification} pays workers
for AI computation. This section makes precise a gap that any such protocol must
close before its rewards mean anything: \emph{a signature proves authorship, a
plurality proves agreement, and neither proves computation.} We state the gap as
a concrete attack against a deployed settlement engine, because it is the reason
the execution-level proofs of Sections~\ref{sec:four-proofs}--\ref{sec:tiers}
exist.

\subsection{The settlement engine, as deployed}

The Lux A-Chain (\texttt{aivm}) settles AI tasks with a commit--reveal quorum.
For a task $t$ naming a model specification $\mathsf{modelSpec}$ and a prompt,
each selected operator $o$ runs (nominally) the model, obtains an output, and
publishes a commitment
\begin{equation}
\label{eq:aivm-commit}
c_o \;=\; \mathsf{keccak}\bigl(t \,\|\, \mathsf{modelSpecHash} \,\|\,
   \mathsf{promptHash} \,\|\, \mathsf{outputHash} \,\|\,
   \mathsf{embeddingHash} \,\|\, o \,\|\, \mathsf{nonce}\bigr),
\end{equation}
then later reveals the preimage. Reveal verification recomputes $c_o$ from the
revealed fields and the caller's address and checks equality; a mismatch is
rejected. Settlement tallies the revealed $\mathsf{outputHash}$ values, takes the
\emph{plurality} group (the statistical mode), and---if that group's size meets a
threshold $\theta$---declares its hash canonical and pays \texttt{rewardPerOperator}
to every operator in it from escrow, slashing those who withheld a reveal.

Two properties of this design are worth quoting from the implementation itself,
because they are stated there without euphemism. The settlement substrate's
engine documentation records that
\begin{quote}\itshape
``\ldots the settlement \textsc{result}---did $\ge$ threshold staked operators
independently submit the same \texttt{output\_hash} under the same
\textsc{ModelSpec}---is what A-Chain consensus agrees on; validators never run the
model.''
\end{quote}
and the address binding in~\eqref{eq:aivm-commit} is documented as
\begin{quote}\itshape
``the anti-copy / anti-front-run control: a peer who observes operator A's commit
cannot replay it as their own \ldots.''
\end{quote}
The first sentence is the safety premise of the entire mechanism; the second names
exactly what the commitment proves. We take both at face value.

\subsection{What is---and is not---proven}

The commitment~\eqref{eq:aivm-commit} binds an output hash to an \emph{author}
(the address $o$ is inside the keccak preimage and is re-checked on reveal) and,
because $\mathsf{nonce}$ and $o$ are mixed in, resists copying and front-running.
The plurality rule binds the canonical answer to \emph{agreement} (how many
distinct authors published the same hash). Neither operation references the model
weights, the prompt tokens beyond their hash, or any artifact of the forward pass.
We isolate the consequence.

\begin{proposition}[Computation is unconstrained by authorship and agreement]
\label{prop:gap}
Let $\Pi$ be a settlement engine whose verdict is a function only of
(i)~address-bound commitments of the form~\eqref{eq:aivm-commit} and
(ii)~the multiset of revealed output hashes via a plurality rule with threshold
$\theta$. Then for any prompt and any model specification, a coalition of
$\theta$ operators that never evaluates the model can cause $\Pi$ to settle---and
to pay each coalition member---on an output hash $h^\star$ of the coalition's
choosing, provided $h^\star$ is the plurality among revealed hashes.
\end{proposition}

\begin{proof}
Each coalition member $o_i$ ($1 \le i \le \theta$) chooses the same target hash
$h^\star$, samples a fresh $\mathsf{nonce}_i$, and computes
$c_{o_i}$ by~\eqref{eq:aivm-commit} with $\mathsf{outputHash} := h^\star$ and its
own address $o_i$. Each $c_{o_i}$ is a valid, distinct, non-replayable commitment:
it is well-formed (a keccak image), it is bound to the correct author (so reveal
verification's recompute-and-compare succeeds), and it differs from every other
member's commitment (distinct addresses and nonces), so the anti-copy control
raises no objection. On reveal, all $\theta$ members publish $h^\star$; their
group has size $\theta \ge \theta$ and, by hypothesis, is the plurality. $\Pi$
declares $h^\star$ canonical and pays each member. No member evaluated the model;
$h^\star$ was never required to be the model's output on the prompt. \qed
\end{proof}

Proposition~\ref{prop:gap} is not exotic. In the degenerate but common case of a
\emph{decision} task---a yes/no classification, a moderation verdict, an
``approve/reject'' agentic action---the output space is one bit, so $h^\star$ is
one of two hashes. The cost-minimizing strategy is then explicit and requires
\emph{no model at all}:

\begin{corollary}[The one-bit guess]
\label{cor:onebit}
For a binary-output task with prior $\Pr[\text{answer}=1]=q\ge\tfrac12$, the
strategy ``commit and reveal the hash of the modal answer $1$, never run the
model'' settles correctly with probability $q$ at zero inference cost, and---when
a coalition of size $\theta$ adopts it---settles \emph{and mints} with probability
$1$ regardless of the true answer.
\end{corollary}

The corollary is the headline: under authorship-plus-agreement, the network can be
made to mint for a one-bit guess. A signature scheme cannot fix this, because the
defect is not forgery---every signature and every commitment in the attack is
genuine. The missing predicate is \emph{``this output is the result of evaluating
the named model on the named input,''} and that predicate is exactly what an
execution-level proof supplies.

\subsection{Why the existing fallbacks do not close it}

Three mechanisms in the deployed stack might be hoped to help; none does, for a
structural reason worth stating.

\begin{itemize}[leftmargin=1.4em]
\item \textbf{Re-execution by validators} is the natural fix, and the A-Chain
documentation contemplates an optional ``$M$-of-$N$ redundant'' mode in which
validators re-run inference. But this is precisely what the deployed engine
\emph{does not} do (``validators never run the model''), and naive re-execution
costs $M\times$ the compute per task, which is why it was not made the default.
Section~\ref{sec:freivalds} replaces $O(n^3)$ re-execution with an $O(n^2)$
\emph{check}, making the fix affordable.
\item \textbf{TEE attestation} (the path of Section~\ref{sec:verification}) does
bind output to enclave execution, and closes the gap \emph{when present}. But it
requires confidential-compute hardware, excludes CPU-only and consumer-GPU
providers, and reduces the trust base to a vendor's remote-attestation root. It is
a strong path, not an egalitarian one (Section~\ref{sec:tiers}).
\item \textbf{A GPU ``proof'' kernel} exists in the A-Chain backend, but it is a
TEE-attestation envelope check---it verifies a signature over a hardware
measurement---and it is never invoked by the settlement or mint path. It is, by
construction, subject to Proposition~\ref{prop:gap}: it proves authorship of an
attestation, not execution of a model.
\end{itemize}

The remainder of this paper \emph{constructs, proves sound, and builds} the missing
predicate as a \emph{software-only} artifact---no enclave, no special hardware---that is a
side effect of the genuine computation, that cannot be produced without performing the
matmuls, and that any participant can both emit and check against a network-shared
difficulty. The predicate is no longer only on paper: the verifier, its emission from a
real forward/graph pass, the off-chain challenger, and the on-chain fail-closed mint gate
are implemented and tested ($38$ Rust, $23$ Go, $87$ Solidity tests;
Section~\ref{sec:impl-status}). So the attack above describes the settlement of a system
that does \emph{not} install the compute-proof gate---which the deployed contracts now
refuse to be: \texttt{AICoinMiner} reverts with \texttt{ProofGateUnavailable} when no
verifier is wired, so a proofless mint is impossible out of the box. The residual is the
floating-point tier and wider architecture emission, not the closing of this attack.
